\documentclass[11pt]{article}
\usepackage[a4paper,margin=1in]{geometry}
\usepackage{fontspec}
\usepackage[boldfont=false]{xeCJK}
\setCJKmainfont{FandolSong-Regular.otf}
\usepackage{amsmath}
\usepackage{amssymb}
\usepackage{array}
\usepackage{booktabs}
\usepackage{graphicx}
\usepackage{adjustbox}
\usepackage{enumitem}
\setlist[itemize]{topsep=2pt,partopsep=0pt,parsep=0pt,itemsep=2pt,leftmargin=1.4em}
\setlist[enumerate]{topsep=2pt,partopsep=0pt,parsep=0pt,itemsep=2pt,leftmargin=1.6em}
\usepackage{parskip}
\usepackage{fvextra}
\fvset{breaklines=true,breakanywhere=true,fontsize=\small}
\setlength{\parindent}{0pt}

\begin{document}

\begin{center}
  {\LARGE\bfseries Electric fields}\\[0.35em]
  {\large A Level Physics}
\end{center}
\vspace{0.6em}

\section*{Electric fields}

An \textbf{electric field} 电场 is a region where a charge feels a \textbf{force} 力 from other charges. The \textbf{electric field strength} 电场强度 $E$ at a point is the \textbf{force per unit positive charge} on a small positive \textbf{test charge} 检验电荷 placed there:

\[
E = \frac{F}{q}.
\]

Unit: $\text{N C}^{-1}$ (the same as $\text{V m}^{-1}$, as we will see). $E$ is a \textbf{vector} 矢量, pointing the way the force acts on a \textbf{positive} charge. The force on a charge $q$ is

\[
F = qE,
\]

opposite to the field if $q$ is negative.

\subsection*{Field lines}

\begin{itemize}
\item \textbf{field lines} 场线 point the way the force acts on a positive test charge.

\item lines start on \textbf{positive} charges and end on \textbf{negative} charges (or go to infinity).

\item lines never cross; closer lines mean a stronger field.

\end{itemize}

Examples: a positive \textbf{point charge} 点电荷 has radial lines pointing out; a negative one has lines pointing in; two opposite charges (a \textbf{dipole} 偶极子) have lines curving from + to −; two parallel charged plates give a \textbf{uniform field} 匀强场 of equally spaced parallel lines.

\par\medskip
\noindent\begin{minipage}{\linewidth}
\centering
\adjustbox{max width=\linewidth, max totalheight=0.42\textheight}{%
  \includegraphics{/Users/joschka/Desktop/Programming/eXercise/output/eXam/handouts/a_level_physics/assets/18-p304-electric-field-patterns.pdf}}\\[0.3em]
\emph{Field-line patterns for parallel plates, a dipole, a point charge, and a charged sphere above an earthed plate}
\end{minipage}\par\medskip

\section*{Uniform electric fields}

Between two parallel plates a distance $d$ apart with \textbf{potential difference} 电势差 $V$ between them, the field is uniform (apart from edge effects) with size

\[
E = \frac{V}{d}.
\]

It points from the higher-potential plate to the lower one. The unit $\text{V m}^{-1}$ comes straight from this and equals $\text{N C}^{-1}$.

\subsection*{A charged particle in a uniform field}

A charge $q$ in a uniform field feels a \textbf{constant force} $F = qE$, so a \textbf{constant acceleration} $a = qE/m$ — just like a mass in a uniform gravitational field.

\begin{itemize}
\item released \textbf{at rest}, it speeds up along the field (positive charge) or against it (negative charge), gaining \textbf{kinetic energy} 动能.

\item entering \textbf{at right angles} to the field, it follows a \textbf{parabolic} 抛物线 path — like a \textbf{projectile} 抛体 in gravity. This is how a \textbf{cathode-ray tube} 阴极射线管 used to steer its beam.

\end{itemize}

\par\medskip
\noindent\begin{minipage}{\linewidth}
\centering
\adjustbox{max width=\linewidth, max totalheight=0.42\textheight}{%
  \includegraphics{/Users/joschka/Desktop/Programming/eXercise/output/eXam/handouts/a_level_physics/assets/18-p309-charged-particle-parabola.pdf}}\\[0.3em]
\emph{A charge entering a uniform field at right angles follows a parabolic path, like a projectile}
\end{minipage}\par\medskip

\section*{Coulomb's law}

For two point charges $Q_{1}$ and $Q_{2}$ a distance $r$ apart in free space, each feels a force of size

\[
F = \frac{Q_{1} Q_{2}}{4\pi\varepsilon_{0} r^{2}}.
\]

This is \textbf{Coulomb's law} 库仑定律. Here $\varepsilon_{0} = 8.85 \times 10^{-12}\ \text{F m}^{-1}$ is the \textbf{permittivity of free space} 真空电容率, and $1/(4\pi\varepsilon_{0}) \approx 9.0 \times 10^{9}\ \text{N m}^{2}\ \text{C}^{-2}$. The force is along the line joining the charges: \textbf{repulsive} for like charges, \textbf{attractive} for opposite charges.

\subsection*{Spheres treated as point charges}

A spherical \textbf{conductor} 导体 with total charge $Q$ gives, at any point \textbf{outside}, the same field as a point charge $Q$ at its centre (measure $r$ from the centre). \textbf{Inside} a hollow charged conductor the field is zero, so the conductor is an \textbf{equipotential} 等势面.

\section*{Electric field due to a point charge}

The field at distance $r$ from a point charge $Q$ is

\[
E = \frac{Q}{4\pi\varepsilon_{0} r^{2}}.
\]

It points out from a positive $Q$, in towards a negative $Q$, and falls as $1/r^{2}$ — just like \textbf{gravitational field} 重力场 strength, except gravity is always attractive. For several charges, add the fields as a \textbf{vector sum} 矢量和.

\section*{Electric potential}

\textbf{Electric potential} 电势 $V$ at a point is the \textbf{work done per unit positive charge} in bringing a small positive test charge from \textbf{infinity} 无穷远 to that point:

\[
V = \frac{W}{q}.
\]

Unit: $\text{V}$. The potential is \textbf{zero at infinity}. For a positive source charge $V > 0$ everywhere outside; for a negative source charge $V < 0$.

\subsection*{Potential due to a point charge}

\[
V = \frac{Q}{4\pi\varepsilon_{0} r}.
\]

Note the $1/r$ here (compared with $1/r^{2}$ for the field). $V$ is a \textbf{scalar} 标量; for several charges, add the potentials (with sign).

\par\medskip
\noindent\begin{minipage}{\linewidth}
\centering
\adjustbox{max width=\linewidth, max totalheight=0.42\textheight}{%
  \includegraphics{/Users/joschka/Desktop/Programming/eXercise/output/eXam/handouts/a_level_physics/assets/18-p313-point-charge-potential.pdf}}\\[0.3em]
\emph{The potential near a point charge varies as 1/r — positive for a positive charge, negative for a negative one}
\end{minipage}\par\medskip

\subsection*{Link between field and potential}

The field equals the \textbf{negative potential gradient} 电势梯度:

\[
E = -\frac{dV}{dx}.
\]

Between parallel plates $V$ changes evenly with position, giving $E = V/d$ as before. The minus sign means the field points towards \textbf{lower} potential. For a point charge, $-\dfrac{dV}{dr} = \dfrac{Q}{4\pi\varepsilon_{0} r^{2}} = E$.

\par\medskip
\noindent\begin{minipage}{\linewidth}
\centering
\adjustbox{max width=\linewidth, max totalheight=0.42\textheight}{%
  \includegraphics{/Users/joschka/Desktop/Programming/eXercise/output/eXam/handouts/a_level_physics/assets/18-p314-potential-field-link.pdf}}\\[0.3em]
\emph{In a uniform field the potential falls steadily with distance, so the field strength V/d is constant}
\end{minipage}\par\medskip

\subsection*{Electric potential energy}

A charge $q$ at a point of potential $V$ has \textbf{electric potential energy} 电势能 $E_{\text{P}} = qV$. For two point charges $Q$ and $q$ a distance $r$ apart:

\[
E_{\text{P}} = \frac{Qq}{4\pi\varepsilon_{0} r}.
\]

\begin{itemize}
\item \textbf{like} charges: $E_{\text{P}} > 0$ — stored energy that would be released if they flew apart.

\item \textbf{opposite} charges: $E_{\text{P}} < 0$ — a \textbf{bound} 束缚 system; energy must be supplied to separate them.

\end{itemize}

In both cases $E_{\text{P}} \to 0$ as $r \to \infty$.

\subsection*{Worked-example pattern}

An \textbf{electron} 电子 orbits a \textbf{nucleus} 原子核 of charge $+Ze$ at distance $r$. The Coulomb attraction provides the \textbf{centripetal force} 向心力:

\[
\frac{Z e^{2}}{4\pi\varepsilon_{0} r^{2}} = \frac{m_{e} v^{2}}{r}, \qquad v = \sqrt{\frac{Z e^{2}}{4\pi\varepsilon_{0} m_{e} r}}.
\]

The total energy is kinetic plus potential:

\[
E_{\text{total}} = \tfrac{1}{2} m_{e} v^{2} - \frac{Z e^{2}}{4\pi\varepsilon_{0} r} = -\frac{Z e^{2}}{8\pi\varepsilon_{0} r},
\]

which is negative (a bound state).

\subsection*{Gravitational versus electric}

The two field theories look alike:

\begin{center}
\adjustbox{max width=\linewidth}{%
\begin{tabular}{lll}
\toprule
\textbf{Quantity} & \textbf{Gravitational} & \textbf{Electric} \\
\midrule
Source & mass $M$ (always positive) & charge $Q$ (can be ±) \\
Field strength & $g = GM/r^{2}$ & $E = Q/(4\pi\varepsilon_{0} r^{2})$ \\
Force on test object & $F = mg$ & $F = qE$ \\
Potential & $\phi = -GM/r$ & $V = Q/(4\pi\varepsilon_{0} r)$ \\
PE of two & $-GMm/r$ & $Qq/(4\pi\varepsilon_{0} r)$ \\
Nature & always attractive & attractive or repulsive \\
\bottomrule
\end{tabular}}
\end{center}

The minus sign in the \textbf{gravitational potential} 引力势 reflects that gravity is always attractive; the electric potential takes the sign of the source charge.


\end{document}
